\subsubsection{Orthogonale Projektion}
\label{subsubsec:vector_algebra_orthogonal_projection}

Mit dem Skalarprodukt lässt sich die orthogonale Projektion
\EnglishTerm{orthogonal projection} eines Vektors auf einen anderen
berechnen. Dies hat im Maschinenbau vielfältige Anwendungen: Wirkt eine Kraft
auf einen Stab, so erzeugt ihre parallele Komponente eine axiale Zug- oder
Druckbeanspruchung und ihre senkrechte Komponente eine Biege- beziehungsweise
Drehwirkung.

Um diese Wirkungen getrennt untersuchen zu können, zerlegen wir die Kraft in
ihre parallele und senkrechte Komponente. Allgemein lässt sich ein Vektor
$\vec b$ relativ zu einer Richtung $\vec a\ne\vec 0$ schreiben als
\[
  \vec b=\vec b_{\parallel}+\vec b_{\perp}.
\]

\begin{center}
  \CachedTikzFigure{content/resources/chapters/vector_algebra/tikz/orthogonal_projection.tex}
\end{center}

Mit dem Einheitsvektor $\vec e_a=\vec a/|\vec a|$ ergibt sich die
vorzeichenbehaftete Komponente von $\vec b$ in Richtung $\vec a$:
\begin{ClozeArea}[34]{5}
  \[
    \begin{aligned}
      \vec b\cdot\vec e_a
      &=(\vec b_{\parallel}+\vec b_{\perp})\cdot\vec e_a\\[0.5em]
      &=\vec b_{\parallel}\cdot\vec e_a
        +\underbrace{\vec b_{\perp}\cdot\vec e_a}_{=0}.
    \end{aligned}
  \]
\end{ClozeArea}

\begin{KeyPointBox}{Parallele und senkrechte Komponente}{vector_algebra_projection_components}
  Die Komponente eines Vektors $\vec{b}$ \dots
  \begin{itemize}
    \item parallel zum senkrecht, zum 
      \[
        \vec b_{\parallel} = (\vec b\cdot\vec e_a)\,\vec e_a
      \]
      %
      \item senkrecht zum ????
        \[
          \vec b_{\perp} = \vec b-\vec b_{\parallel}.
        \]
  \end{itemize} 


  Dabei gilt $|\vec b_{\parallel}|=|\vec b\cdot\vec e_a|$ und
  $\vec a\cdot\vec b_{\perp}=0$.
\end{KeyPointBox}

\begin{ExampleBox}{Einen Vektor orthogonal zerlegen}{vector_algebra_projection_example}
  \noindent
  \begin{minipage}[c]{0.57\linewidth}
    Gegeben seien
    \[
      \textcolor{THMBlue}{\vec a}=\begin{pmatrix}4\\2\end{pmatrix},
      \qquad
      \textcolor{olive!80!black}{\vec b}=\begin{pmatrix}1\\3\end{pmatrix}.
    \]
    Dann gilt
    \[
      \begin{aligned}
        \vec b_{\parallel}
        &=\frac{\textcolor{THMBlue}{\vec a}\cdot
            \textcolor{olive!80!black}{\vec b}}
          {\textcolor{THMBlue}{\vec a}\cdot
            \textcolor{THMBlue}{\vec a}}\,
          \textcolor{THMBlue}{\vec a}
          =\frac{10}{20}\begin{pmatrix}4\\2\end{pmatrix}\\
        &=\Cloze{2.5cm}{\begin{pmatrix}2\\1\end{pmatrix}},
      \end{aligned}
    \]
    \[
      \textcolor{black!60}{\vec b_{\perp}}
      =\textcolor{olive!80!black}{\vec b}-\vec b_{\parallel}
      =\Cloze{2.5cm}{\begin{pmatrix}-1\\2\end{pmatrix}},
      \qquad
      \textcolor{THMBlue}{\vec a}\cdot
      \textcolor{black!60}{\vec b_{\perp}}=0.
    \]
  \end{minipage}\hfill
  \begin{minipage}[c]{0.39\linewidth}
    \centering
    \begin{tikzpicture}[>=Latex,x=0.72cm,y=0.72cm]
      \draw[step=1cm,black!10,line width=0.35pt]
        (-1.35,-0.45) grid (4.45,3.45);
      \draw[->,black!55,line width=0.75pt]
        (-1.35,0)--(4.55,0) node[right] {$x_1$};
      \draw[->,black!55,line width=0.75pt]
        (0,-0.45)--(0,3.55) node[above] {$x_2$};
      \foreach \x in {-1,1,2,3,4}
        \draw[black!55] (\x,0.07)--(\x,-0.07)
          node[below=2pt] {$\x$};
      \foreach \y in {1,2,3}
        \draw[black!55] (0.07,\y)--(-0.07,\y)
          node[left=2pt] {$\y$};

      \coordinate (O) at (0,0);
      \coordinate (A) at (4,2);
      \coordinate (P) at (2,1);
      \coordinate (B) at (1,3);
      \draw[->,THMBlue,line width=1.35pt] (O)--(A)
        node[pos=0.82,below right] {$\vec a$};
      \draw[->,olive!80!black,line width=1.35pt] (O)--(B)
        node[pos=0.72,left] {$\vec b$};

      \IfClozeTF{}{%
        \draw[->,black,line width=1.25pt] (O)--(P)
          node[pos=0.58,below right] {$\vec b_{\parallel}$};
        \draw[->,black!60,line width=1.25pt] (P)--(B)
          node[midway,right] {$\vec b_{\perp}$};
        \draw[black!55,line width=0.75pt]
          ($(P)!0.16!(O)$)--++(-0.16,0.32)--++(0.32,0.16);
      }%
      \fill[black] (O) circle[radius=1.2pt];
    \end{tikzpicture}
  \end{minipage}
\end{ExampleBox}
